\section{Forcing activity and excluding full periodicity}
\label{sec:activation}

The exclusion equations constrain which labels can be active together.
We now add one finite tile for each channel to ensure that every ordinary
channel is active at every point, and that each seed is active somewhere.
Throughout this section, we retain the groups and coordinates of
\cref{sec:encoding}.  In particular, the useful output of channel $i$ is
$c_i\in P_i=\Zmod{a_i}\times\Zmod{b_i}$, its low input is $k_i\in K_i$,
and its sheared high coordinate is $\beta_i\in K_i$.

\subsection{A shift list with two fibre orderings}

The shift list will serve two purposes: its nonuniform multiplicities
will force activity, while its uniform counts in each coordinate will
permit the explicit common solution in \cref{sec:model}.

Fix a channel $i$.  For $e\in P_i$, let $\mathbf 1_{\{e\}}$ denote the
indicator of the singleton $\{e\}$, and define the integer-valued functions
\[
 \omega_i
 =\mathbf 1_{\{(0,0)\}}-\mathbf 1_{\{(1,0)\}}
  -\mathbf 1_{\{(0,1)\}}+\mathbf 1_{\{(1,1)\}},
 \qquad
 \mu_i=1+\omega_i
 \quad\text{on }P_i.
\]
Here $1$ denotes the constant function.  The function $\mu_i$ has values
two at $(0,0)$ and $(1,1)$, zero at $(1,0)$ and $(0,1)$, and one elsewhere.
It is therefore a nonnegative, nonconstant multiplicity function.
Its total mass is $a_i b_i$, and its marginals are uniform:
\[
 \sum_{\alpha\in\Zmod{a_i}}\mu_i(\alpha,\xi)=a_i
 \quad(\xi\in\Zmod{b_i}),
 \qquad
 \sum_{\xi\in\Zmod{b_i}}\mu_i(\alpha,\xi)=b_i
 \quad(\alpha\in\Zmod{a_i}).
\]
Choose an index set $\Delta_i$ of size $a_i b_i$ and a list
$e_\delta=(e_\delta^a,e_\delta^b)\in P_i$, $\delta\in\Delta_i$, in
which each $e\in P_i$ occurs $\mu_i(e)$ times.  The following notation
for the list and its orderings is local to channel $i$.

For each $\xi$, order the $a_i$ indices with $e_\delta^b=\xi$ by an
arbitrary bijection to $\Zmod{a_i}$, and define $\rho^a(\delta)$ by
subtracting $e_\delta^a$ from that assigned value.  Independently, for
each $\alpha$, order the $b_i$ indices with $e_\delta^a=\alpha$ by a
bijection to $\Zmod{b_i}$, and subtract $e_\delta^b$ to define
$\rho^b(\delta)$.  Thus both maps in
\begin{equation}
\label{act:fibre-orderings}
\begin{aligned}
 \{\delta\in\Delta_i:e_\delta^b=\xi\}
   &\longrightarrow\Zmod{a_i},
 &\delta&\longmapsto e_\delta^a+\rho^a(\delta),\\
 \{\delta\in\Delta_i:e_\delta^a=\alpha\}
   &\longrightarrow\Zmod{b_i},
 &\delta&\longmapsto e_\delta^b+\rho^b(\delta)
\end{aligned}
\end{equation}
are bijections.  Repeated values of $e_\delta$ have distinct indices;
the orderings act on these indices.

\subsection{The activation tiles}

For a channel $\ell$, write
$R_\ell=r_\ell\Zmod{r_\ell^2}$ for the subgroup of offsets with low
coordinate zero.  It has $r_\ell$ elements.  Given $h\in\Z^2$,
$e\in P_i$, and $E\subseteq\Zmod{D}$, define a batch of offsets by
\[
\begin{split}
 \mathcal B_i(h,e;E)=
 \bigl\{\bigl(h,\zeta,(w_\ell,d_\ell)_{\ell\in\mathcal I}\bigr):
 &\ \zeta\in E,\quad
 w_i=h_1\pmod{r_i^2},\quad d_i=e,\\[-2pt]
 &\ w_\ell\in R_\ell,\quad d_\ell\in P_\ell
       \text{ for every }\ell\ne i\bigr\}.
\end{split}
\]
Every combination of the freely varied coordinates occurs once in this
set.

For an ordinary channel $i\in\{0,\ldots,N-1\}$ and each pair
$(l,\delta)\in K_i\times\Delta_i$, choose an integer multiple
$h=h^{(i)}_{l,\delta}$ of $u_i=(1,-i)$ such that
\begin{equation}
\label{act:ordinary-crt}
 h_1\equiv l\pmod{r_i},\qquad
 h_1\equiv\rho^a(\delta)\pmod{a_i},\qquad
 h_1\equiv\rho^b(\delta)\pmod{b_i}.
\end{equation}
The moduli are pairwise coprime, so the Chinese remainder theorem gives
infinitely many choices of $h_1$.  Choose these horizontal offsets
distinctly for the different pairs $(l,\delta)$, and set
\[
 F_i^{\mathrm{act}}
 =\bigcup_{(l,\delta)\in K_i\times\Delta_i}
    \mathcal B_i\bigl(h^{(i)}_{l,\delta},e_\delta;\Zmod{D}\bigr).
\]
The restriction $h\in\Z u_i$ gives $L_i(x-h)=L_i(x)$ for every $x$.

For a seed $i\in\{\eta_1,\eta_2\}$, let $i'$ denote the other seed.
For every $(l,\delta,\tau)\in K_i\times\Delta_i\times S$, choose
$h=h^{(i)}_{l,\delta,\tau}\in\Z^2$ satisfying
\cref{act:ordinary-crt} and the additional conditions
\begin{equation}
\label{act:seed-crt}
 s(h)=\tau,\qquad h_1\equiv0\pmod{a_{i'}}.
\end{equation}
There is no slope restriction in this case.  Writing
$\tau=(\tau_1,\tau_2)\in(\Zmod{p^2})^2$, the first coordinate is
prescribed modulo the pairwise coprime integers
$r_i,a_i,b_i,p^2,a_{i'}$, while the second coordinate only needs to
satisfy $h_2\equiv\tau_2\pmod{p^2}$.  Again there are infinitely many
choices, so take all horizontal offsets within this seed tile to be
distinct.  Define
\[
 F_i^{\mathrm{act}}
 =\bigcup_{(l,\delta,\tau)\in K_i\times\Delta_i\times S}
    \mathcal B_i\bigl(h^{(i)}_{l,\delta,\tau},e_\delta;\{0\}\bigr).
\]
Thus an ordinary batch varies the $z$-offset freely, whereas a seed batch
has $z$-offset zero.  The congruence involving $a_{i'}$ will allow both
seed tests to hold with one common $z$-output in the construction of
\cref{model:system}.

All unions just defined are disjoint unions of finite sets, because
different batches have different horizontal offsets.  In particular,
the repeated entries of the auxiliary shift list do not introduce
weighted tile elements.  These tiles are nonempty, and direct counting
gives, for every channel $i$,
\[
 \abs{F_i^{\mathrm{act}}}
   =D\prod_{\ell\in\mathcal I}r_\ell\abs{P_\ell}
   =\abs{H_0}.
\]
We add the equations $A\oplus F_i^{\mathrm{act}}=G$ to the graph,
dependence, and exclusion equations.

\subsection{Exact coverage in a fibre}

\begin{lemma}[Activation fibre criterion]
\label{act:fibre-criterion}
Suppose $A\oplus H_0=G$, so that $A$ is a graph over $(x,k)$.
For a target base point $(x,k)$ and a batch with horizontal offset $h$,
write
\begin{equation}
\label{act:source-base}
 x'=x-h,\qquad k_i'=k_i-l,\qquad k_\ell'=k_\ell\quad(\ell\ne i).
\end{equation}
For an ordinary channel $i$, the equation
$A\oplus F_i^{\mathrm{act}}=G$ holds if and only if, at every target
base point, the map
\begin{equation}
\label{act:ordinary-map}
 (l,\delta)\longmapsto
 \bigl(c_i(x',k')+e_\delta,\,\beta_i(x',k')\bigr)
\end{equation}
is a bijection from $K_i\times\Delta_i$ to $P_i\times K_i$.
For a seed $i$, the corresponding necessary and sufficient condition is
that, at every target base point,
\begin{equation}
\label{act:seed-map}
 (l,\delta,\tau)\longmapsto
 \bigl(c_i(x',k')+e_\delta,\,\beta_i(x',k'),\,z(x',k')\bigr)
\end{equation}
be a bijection from $K_i\times\Delta_i\times S$ to
$P_i\times K_i\times\Zmod{D}$.
\end{lemma}

\begin{proof}
Fix a target base point and a batch.  Every offset in that batch has the
same horizontal and low coordinates, so the graph equation supplies
exactly one possible source point of $A$: the point above the base in
\eqref{act:source-base}.  The $i$th low-coordinate shift is $l$ because
$h_1\equiv l\pmod{r_i}$.  Since $x'_1=x_1-h_1$, we have
\[
 [k_i'-x'_1]_{r_i}=[k_i-x_1]_{r_i}.
\]
Consequently the shear formula \eqref{enc:shear} gives
\[
 v_i(x',k')+h_1
   =x_1+[k_i-x_1]_{r_i}+r_i\beta_i(x',k')
   \quad\text{in }\Zmod{r_i^2}.
\]
Thus the offset preserves the sheared high coordinate $\beta_i$ exactly,
and the target useful output is $c_i(x',k')+e_\delta$.

For every $\ell\ne i$, adding the freely chosen element of $R_\ell$
fills the entire high-coordinate fibre with low coordinate $k_\ell$
once.  Independently, adding the freely chosen element of $P_\ell$
fills that useful-output coordinate once.  In an ordinary batch the
$z$-coordinate is also filled once; in a seed batch it remains equal
to $z(x',k')$.  Hence each batch covers exactly the part of the target
fibre specified by its tuple in \cref{act:ordinary-map} or
\cref{act:seed-map}, with one representation for every choice of the
remaining coordinates.  Distinct batches correspond to distinct tile
elements, so exact coverage by the whole tile is equivalent to every
specified tuple occurring once.  This is precisely the asserted
bijection condition.
\end{proof}

In particular, exact coverage forces every value of the target $c_i$
to occur $r_i$ times among the ordinary batches, or $D r_i$ times
among the seed batches.  These necessary marginal counts already force
the activity we need.

\subsection{Every required channel is active}

\begin{lemma}[Ordinary activation]
\label{act:ordinary}
Every ordinary channel has at least one active label at every
$x\in\Z^2$ in any common solution of the equations specified so far.
\end{lemma}

\begin{proof}
Suppose channel $i$ has no active label at some $x$.  A function on the
finite product $K_i=\prod_{j\in J_i}\Zmod{r_{ij}}$ that is independent
of every digit is constant: any two inputs can be joined by changing
one digit at a time.  Thus $g_{i,x}$ has a constant value $c_0\in P_i$.
By the dependence condition \eqref{enc:dependence} and the equality
$L_i(x-h)=L_i(x)$, every source in the activation test at $(x,k)$ has
$c_i(x',k')=c_0$.  For each $\delta$ there are $r_i$ choices of $l$.
The number of batches giving any target value $c\in P_i$ is therefore
\[
 r_i\mu_i(c-c_0).
\]
The fibre criterion requires this number to be $r_i$ for every $c$.
That is impossible because $\mu_i$ is not constant.
\end{proof}

\begin{lemma}[Seed activation]
\label{act:seed}
Each of the two seeds has its label active at some point of $\Z^2$
in any common solution of the equations specified so far.
\end{lemma}

\begin{proof}
Fix a seed $i$ and suppose it is nowhere active.  The dependence
condition and the absence of activity give a function $q:S\to P_i$
such that $c_i(x,k)=q(s(x))$ for every base point.  Define its
nonnegative integer histogram by
\[
 Q(\alpha,\xi)
 =\abs{\{s\in S:q(s)=(\alpha,\xi)\}},
 \qquad
 \sum_{(\alpha,\xi)\in P_i}Q(\alpha,\xi)=D.
\]
At any target base point and for a fixed $\delta$, the source pairs
\[
 (s(x'),k_i')=(s(x)-\tau,k_i-l)
\]
run once through $S\times K_i$ as $(\tau,l)$ varies.  Thus the
number of seed batches with useful output $c\in P_i$ is
$r_i\sum_{\delta\in\Delta_i}Q(c-e_\delta)$.  The necessary marginal
count from \cref{act:fibre-criterion}, divided by $r_i$, yields
\[
 (Q*\mu_i)(c)=D\qquad(c\in P_i),
\]
where $(Q*\mu_i)(c)=\sum_{e\in P_i}\mu_i(e)Q(c-e)$ is convolution
on the finite group $P_i$.  Convolution with the constant function
$1$ already equals $D$, so $Q*\omega_i=0$.  Written out, this says
\begin{equation}
\label{act:mixed-difference}
 Q(\alpha,\xi)-Q(\alpha-1,\xi)
 -Q(\alpha,\xi-1)+Q(\alpha-1,\xi-1)=0.
\end{equation}

We spell out the integer and positivity consequences of this identity.
For each $\alpha$, the difference
$Q(\alpha,\xi)-Q(\alpha-1,\xi)$ is independent of $\xi$, because
successive differences in the second coordinate vanish.  Summing such
differences along the cyclic first coordinate shows that, for any
$\alpha_0$,
\[
 Q(\alpha,\xi)-Q(\alpha_0,\xi)
   =Q(\alpha,0)-Q(\alpha_0,0).
\]
Choose $\alpha_0$ to minimize $Q(\alpha,0)$, and put
\[
 u(\alpha)=Q(\alpha,0)-Q(\alpha_0,0),
 \qquad v(\xi)=Q(\alpha_0,\xi).
\]
Both functions take nonnegative integer values, and
$Q(\alpha,\xi)=u(\alpha)+v(\xi)$.  Summing gives
\[
 D=b_i\sum_{\alpha\in\Zmod{a_i}}u(\alpha)
     +a_i\sum_{\xi\in\Zmod{b_i}}v(\xi).
\]
Since $b_i>D$ and both sums are nonnegative integers, the first sum
must vanish.  It follows that $a_i$ divides $D=p^4$.  But the two
seed values of $a_i$ are $2$ and $3$, whereas $p>200$ is prime.  This
contradiction proves the claim.
\end{proof}

\subsection{Extraction of the line array}

The graph, dependence, exclusion, and activation equations now form a
finite system of tiling equations in $G$.  Its solvability will be
proved in the next section.  We can already exclude all fully periodic
common solutions.

\begin{proposition}
\label{act:nonperiodic}
No common solution of this finite system is invariant under a
finite-index subgroup of $G$.
\end{proposition}

\begin{proof}
Let $A$ be any common solution.  Fix a total ordering of $\Sigma$.
For each ordinary channel $n$ and integer $m$, define
\[
 W(n,m)=\min\mathcal A_n((0,m)),
\]
where the minimum is taken in that ordering.  This is defined by
\cref{act:ordinary}.  The dependence condition implies that
$\mathcal A_n(x)$ depends only on $L_n(x)$; since $L_n(0,m)=m$, the
selected symbol at any $x$ is exactly $W(n,L_n(x))$.

For any integers $d,e$, consider $x=(d,e)$.  If the word
$(W(n,dn+e))_{0\le n<N}$ were forbidden, all its selected labels
would be active at $x$.  This contradicts the corresponding exclusion
equations, by \cref{enc:exclusion}.  Hence $W$ satisfies the line rule.

By \cref{act:seed}, for each $t\in\{1,2\}$ there is a point
$x^{(t)}$ where the label of $\eta_t$ is active.  The seed-pair
exclusions imply that $\mathcal A_n(x^{(t)})$ contains no symbol other
than $t$.  Since this set is nonempty, it equals $\{t\}$.  Therefore
\[
 W\bigl(n,L_n(x^{(1)})\bigr)=1,
 \qquad
 W\bigl(n,L_n(x^{(2)})\bigr)=2
 \quad\text{for every }0\le n<N.
\]
The two arguments in each column must be distinct, and every column
of $W$ is nonconstant.

Suppose now that $A$ is invariant under a finite-index subgroup
$\Lambda\le G$.  Consider the actual group element
\[
 e_{\mathrm v}=\bigl((0,1),0_V\bigr)\in G.
\]
Its coset has finite order in $G/\Lambda$, so some integer $M>0$
satisfies $M e_{\mathrm v}\in\Lambda$.  We thus obtain a period
with \emph{zero finite-coordinate component}:
\[
 A+\bigl((0,M),0_V\bigr)=A.
\]
Translation by this element takes the unique graph point above
$(x,k)$ to the unique graph point above $(x+(0,M),k)$ and leaves all
useful outputs unchanged.  It follows that every active-label set is
unchanged under $x\mapsto x+(0,M)$.  At $x=(0,m)$ this gives
$W(n,m+M)=W(n,m)$ for all $n,m$.
The array is therefore vertically periodic, satisfies the line rule,
and has every column nonconstant, contradicting
\cref{word:obstruction}.
\end{proof}
